\section{explanation}
This section explains the math behind ECDSA.
\subsection{The Discrete Logarithm Problem}
The rising prevalence of elliptice curve cryptography (ECC) is predominantly caused by the strength-per-key-bit, which is substantially greater in algorithms that use this mathematical concept. In return, the increased stregth-per-key-bit and therefore higher security of ECC, is based on the hardness of the elliptic curve discrete logarithm problem (ECDLP), which is assumed to be computationally intractable. The more general discrete logarithm problem (DLP), which it is derived from is also subject to other well known cryptographical algorithms such as the Diffi e-Hellman key exchange. It is defined as follows: Let $G$ be a cyclic group of order $n$ with the generator $g$. The discrete logarithm of $h \in G$ to the base $g$, denoted by $\log_g(h)$, is the unique integer $k$, $0 \leq k \leq n-1$, such that $[k]g = h$. In more broader terms, this essentially states that given a generator $g$ and an element $h$ in a finite cyclic group $G$, find the integer $k$ such that $g^k = h$.
\subsection{ECDSA Domain Parameters}
In case of ECC the cyclic group from the DLP can be thought of as a set of rational Points on an elliptic curve $E$ over a finite field $\mathbb{F}_p$ with $p$ being prime and $E$ representing the Weierstraß equation
$$y^2=x^3+ax+b, a, b \in \mathbb{F}_p$$
The set of rational points in $E$ over $\mathbb{F}_p$ denoted by $E(\mathbb{F}_p)$ can then be written as 
$$E(\mathbb{F}_p) = \{ (x, y) \in \mathbb{F}_p^2 : y^2 = x^3 + ax + b\} \cup \{O\}$$ 
with $O$ being the point at infinity and representing the groups identity element. Why that is, one has to understand the group operation of elliptic curv groups.
The group operation of points in $E(\mathbb{F}_p)$, commonly denoted as addition, has a simple geometric interpretation: 
\begin{figure}
\centering
\includegraphics[width=.5\textwidth]{ecc-point-addition.jpg}
\caption{title}
\end{figure}
The Generator $G$ is the starting  point on $E$ and generates a cyclic subgroub $\langle G \rangle$ of order $n$ in $E(\mathbb{F}_p)$: 
$$\langle G \rangle = \{[1]G, [2]G, [3]G, ..., [n-1]G, [n]G\}$$ 
$k$ is some large integer that defnies the number of times the group operation is performed on the generator $G$ such that $[k]G=P$ where $P \in \langle G \rangle$.
\subsection{calculating $x_r$ and $y_r$}
To perform the group operation on Points on the curve, the formulae for the $x_r$ and $y_r$ value of the resulting Point $R$ have to be derived. Therefore one has to distinguish between two cases. \newline
1. The two points have different $x$ values\newline
2. The two points have the same $x$ and $y$ values\newline
Given two Points $P \in E(\mathbb{F}_p)$ and $Q \in E(\mathbb{F}_p)$ the values for $x_r$ and $y_r$ of the resulting point $R \in E(\mathbb{F}_p)$ can be calculated as follows:
$$\lambda = \frac{y_2-y_1}{x_2-x_1}, y-y_1=\lambda(x-x_1)$$
$$y=\lambda x - \lambda x_1 + y_1, \beta=y_1-\lambda x_1, y=\lambda x + \beta$$
$$(\lambda x + \beta)^2 = x^3 + ax + b$$
$$0 = x^3 - \lambda^2x^2 - 2\lambda x \beta - \beta^2 + ax + b$$
using Vietas formula for qubic polynomials we get:
$$x_3 = \lambda^2 - x_1 - x_2$$
and for $y$ after plugging in x in the first equation:
$$y_3 = \lambda(x_1 - x_3) - y_1$$
For the case in which the two Points are the same, whe simply calculate the tangent to the curve at the point by taking the devirative of e:
$$\frac{dx}{dy} = \frac{d}{dy}(x^3+ax+b)^\frac{1}{2}$$
$$\frac{dx}{dy} = \frac{3x^2 + a}{2\sqrt{x^3+ax+b}}$$
and since $y^2 = x^3 + ax+b$ we can just substitute $y$ back in:
$$\lambda = \frac{3x^2 + a}{2y}$$
Since in this case the points $P$ and $Q$ are technically the same we get the following for $x$ by plugging in, while $y$ will stay the same:
$$x_R = \lambda^2 - 2x_P, y_R = \lambda(x_P - x_R) - y_P$$
\subsection{Key Pair Generation}
Since the ECDSA resides in the realm of asynchronous cryptography, the cryptographic method is based on key pairs consisting of a public key P and a private key d. In case of ECDSA the goals is to prove the authenticity of the sender. Therefore the sender has a private key which he uses to sign the message to send. He also has a public key, which can be used by the receivers of the message, to verify that the message was signed with the corresponding private key and thus came from the expected sender. In the ECDSA the private key is an interger $d$ randomly chosen from the interval $\{1,2,...,n-1\}$ where $n$ is the group order of the cyclic group $G$. The public key $P$ is then generated by applying the group operation on the generator $G$ exactly $d$ times:
\begin{enumerate}
\item choose private key $d = random(\{1,2,...,n-1\})$
\item calculate public key $P = [d]G$
\item output key pair $(d,P)$
\end{enumerate}
\subsection{Signature Generation}
\begin{enumerate}
\item $k = RNG(\{1,2,...,n-1\})$
\item $Q = [k]G$
\item $r = x_Q mod{n}$
\item $k_{inv} = k^{-1} mod{n}$
\item $s = k_{inv} * (r * d_A + H_r(M)) mod{n}$
\item Output $(r,s)$
\end{enumerate}
\subsection{Signature Verification }
\begin{enumerate}
\item verify that $r, s \in \{1, 2, ..., n -1\}$
\item calculate $s_{inv} = s^{-1} mod n$
\item calculate $u_1=s_{inv} * H_r(M) mod{n}$ and $ u_2 = s_{inv} * r mod{n}$
\item $Q = [u_1]G + [u_2]P_A$
\item $v = x_Q mod{n}$
\item Output true if $v =r$, and false otherwise
\end{enumerate}


 



 
